\documentclass[12pt,a4paper]{article}

% PAQUETES
\usepackage[utf8]{inputenc}
\usepackage[T1]{fontenc}
\usepackage[spanish]{babel}
\usepackage{geometry}
\usepackage{xcolor}
\usepackage{listings}
\usepackage{titlesec}
\usepackage{enumitem}
\usepackage{graphicx}

% CONFIGURACION DE PAGINA
\geometry{margin=2.5cm}

% COLORES PARA CODIGO
\definecolor{codegreen}{rgb}{0,0.5,0}
\definecolor{codegray}{rgb}{0.4,0.4,0.4}
\definecolor{codeblue}{rgb}{0,0,0.7}
\definecolor{backcolour}{rgb}{0.96,0.96,0.96}

% CONFIGURACION DE C++
\lstdefinestyle{cppstyle}{
    language=C++,
    backgroundcolor=\color{backcolour},
    commentstyle=\color{codegreen},
    keywordstyle=\color{codeblue}\bfseries,
    numberstyle=\tiny\color{codegray},
    stringstyle=\color{red},
    basicstyle=\ttfamily\small,
    breaklines=true,
    numbers=left,
    numbersep=8pt,
    showstringspaces=false,
    tabsize=4,
    frame=single
}

\begin{document}

% PORTADA
\begin{center}
    {\Large \textbf{UNIVERSIDAD NACIONAL DEL ALTIPLANO}}\\[0.3cm]
    {\large Facultad de Ingeniería Estadística e Informática}\\[0.2cm]
    {\large Escuela Profesional Ingeniería Estadística e Informática}\\[1cm]

    \rule{\textwidth}{0.5pt}\\[0.4cm]
    {\LARGE \textbf{EJERCICIOS DE PILAS (STACK)}}\\[0.3cm]
    {\large Estructura de Datos}\\[0.4cm]
    \rule{\textwidth}{0.5pt}
\end{center}

\vspace{0.8cm}

\noindent
\textbf{Estudiante:} Alberto Roy Chambilla Murillo

\vspace{0.3cm}
\noindent
\textbf{Docente:} Prof. Fred Torres Cruz

\vspace{0.3cm}
\noindent
\textbf{Fecha:} 30/10/2026

\section*{Introducción}

Una pila (Stack) es una estructura de datos que trabaja bajo
el principio LIFO (Last In, First Out), que significa que el
último elemento en ingresar es el primero en salir.

En C++, se puede implementar utilizando la biblioteca
\texttt{<stack>} y sus operaciones principales:
\texttt{push()}, \texttt{pop()}, \texttt{top()} y
\texttt{empty()}.

\section{Ejercicio 1: Apilar números enteros}

\textbf{Objetivo:} Ingresar cinco números enteros,
almacenarlos en una pila y retirarlos uno por one
hasta dejarla vacía.

\subsection*{Código en C++}

\begin{lstlisting}[style=cppstyle]
#include <iostream>
#include <stack>
using namespace std;

int main() {
    stack<int> pila;
    int numero;

    for (int i = 0; i < 5; i++) {
        cout << "Ingrese un numero: ";
        cin >> numero;
        pila.push(numero);
    }

    cout << "\nElementos desapilados: ";

    while (!pila.empty()) {
        cout << pila.top() << " ";
        pila.pop();
    }

    return 0;
}
\end{lstlisting}

\begin{figure}[htbp]
    \centering
    \includegraphics[width=0.5\linewidth]{1.png}
    \caption{Ejecución del Ejercicio 1}
    \label{fig:ej1}
\end{figure}

\subsection*{Ejemplo de ejecución}

\begin{verbatim}
Entrada: 4, 8, 12, 16, 20
Salida: 20, 16, 12, 8, 4
\end{verbatim}

\textbf{Explicación:} Los números se insertan mediante
\texttt{push()}. Posteriormente, se utiliza \texttt{top()}
para consultar el elemento superior y \texttt{pop()} para
eliminarlo. De esta manera, se obtiene el orden inverso
de inserción.

\section{Ejercicio 2: Invertir una palabra}

\textbf{Objetivo:} Leer una palabra, almacenar cada uno
de sus caracteres en una pila y mostrarla en orden inverso.

\subsection*{Código en C++}

\begin{lstlisting}[style=cppstyle]
#include <iostream>
#include <stack>
#include <string>
using namespace std;

int main() {
    stack<char> pila;
    string palabra;

    cout << "Ingrese una palabra: ";
    cin >> palabra;

    for (char letra : palabra) {
        pila.push(letra);
    }

    cout << "Palabra invertida: ";

    while (!pila.empty()) {
        cout << pila.top();
        pila.pop();
    }

    return 0;
}
\end{lstlisting}

\begin{figure}[htbp]
    \centering
    \includegraphics[width=0.5\linewidth]{2.png}
    \caption{Ejecución del Ejercicio 2}
    \label{fig:ej2}
\end{figure}

\subsection*{Ejemplo de ejecución}

\begin{verbatim}
Entrada: DATOS
Salida: SOTAD
\end{verbatim}

\textbf{Explicación:} La pila permite invertir una palabra
porque el último carácter insertado es el primero en salir,
cumpliendo el principio LIFO.

\section{Ejercicio 3: Verificar paréntesis balanceados}

\textbf{Objetivo:} Determinar si los paréntesis de una
expresión están correctamente balanceados.

\subsection*{Código en C++}

\begin{lstlisting}[style=cppstyle]
#include <iostream>
#include <stack>
#include <string>
using namespace std;

int main() {
    stack<char> pila;
    string expresion;
    bool correcto = true;

    cout << "Ingrese una expresion: ";
    getline(cin, expresion);

    for (char caracter : expresion) {
        if (caracter == '(') {
            pila.push(caracter);
        }
        else if (caracter == ')') {
            if (pila.empty()) {
                correcto = false;
                break;
            }
            pila.pop();
        }
    }

    if (!pila.empty()) {
        correcto = false;
    }

    if (correcto) {
        cout << "Los parentesis estan balanceados.";
    }
    else {
        cout << "Los parentesis NO estan balanceados.";
    }

    return 0;
}
\end{lstlisting}

\begin{figure}[htbp]
    \centering
    \includegraphics[width=0.5\linewidth]{3.png}
    \caption{Ejecución del Ejercicio 3}
    \label{fig:ej3}
\end{figure}

\subsection*{Ejemplos de ejecución}

\begin{verbatim}
Entrada: (a+b)*(c-d)
Salida: Los parentesis estan balanceados.

Entrada: ((a+b)
Salida: Los parentesis NO estan balanceados.
\end{verbatim}

\textbf{Explicación:} Cada paréntesis de apertura se
almacena en la pila. Cuando aparece uno de cierre,
se elimina una apertura pendiente. Si no existe una
apertura disponible o quedan elementos al finalizar,
la expresión no está balanceada.

\section{Ejercicio 4: Historial de acciones (Undo)}

\textbf{Objetivo:} Simular la función deshacer de un editor,
almacenando las acciones realizadas en una pila.

\subsection*{Código en C++}

\begin{lstlisting}[style=cppstyle]
#include <iostream>
#include <stack>
#include <string>
using namespace std;

int main() {
    stack<string> pila;
    int opcion;
    string accion;

    do {
        cout << "\n--- HISTORIAL DE ACCIONES ---\n";
        cout << "1. Registrar accion\n";
        cout << "2. Deshacer accion\n";
        cout << "3. Mostrar accion en el tope\n";
        cout << "4. Salir\n";
        cout << "Seleccione una opcion: ";
        cin >> opcion;
        cin.ignore();

        switch (opcion) {
            case 1:
                cout << "Ingrese la accion: ";
                getline(cin, accion);
                pila.push(accion);
                cout << "Accion registrada.\n";
                break;

            case 2:
                if (!pila.empty()) {
                    cout << "Deshaciendo: "
                         << pila.top() << endl;
                    pila.pop();
                }
                else {
                    cout << "No hay acciones para deshacer.\n";
                }
                break;

            case 3:
                if (!pila.empty()) {
                    cout << "Accion en el tope: "
                         << pila.top() << endl;
                }
                else {
                    cout << "El historial esta vacio.\n";
                }
                break;

            case 4:
                cout << "Saliendo del programa.\n";
                break;

            default:
                cout << "Opcion no valida.\n";
        }

    } while (opcion != 4);

    return 0;
}
\end{lstlisting}

\begin{figure}[htbp]
    \centering
    \includegraphics[width=0.5\linewidth]{4.png}
    \caption{Ejecución del Ejercicio 4}
    \label{fig:ej4}
\end{figure}

\subsection*{Ejemplo de ejecución}

\begin{verbatim}
Registrar: Escribir texto
Registrar: Copiar
Registrar: Eliminar

Deshacer: Eliminar
\end{verbatim}

\textbf{Explicación:} Cada acción se almacena en la pila.
Al seleccionar la opción deshacer, se muestra la última
acción mediante \texttt{top()} y se elimina con
\texttt{pop()}.

\section{Ejercicio 5: Conversión de decimal a binario}

\textbf{Objetivo:} Convertir un número decimal positivo
a binario utilizando una pila para almacenar los residuos
de las divisiones sucesivas entre dos.

\subsection*{Código en C++}

\begin{lstlisting}[style=cppstyle]
#include <iostream>
#include <stack>
using namespace std;

int main() {
    stack<int> pila;
    int numero;

    cout << "Ingrese un numero decimal positivo: ";
    cin >> numero;

    if (numero < 0) {
        cout << "Ingrese un numero positivo.";
        return 0;
    }

    if (numero == 0) {
        cout << "Binario: 0";
        return 0;
    }

    while (numero > 0) {
        pila.push(numero % 2);
        numero = numero / 2;
    }

    cout << "Numero en binario: ";

    while (!pila.empty()) {
        cout << pila.top();
        pila.pop();
    }

    return 0;
}
\end{lstlisting}

\begin{figure}[htbp]
    \centering
    \includegraphics[width=0.5\linewidth]{5.png}
    \caption{Ejecución del Ejercicio 5}
    \label{fig:ej5}
\end{figure}

\subsection*{Ejemplo de conversión}

Para convertir el número decimal 13 a binario:

\begin{center}
\begin{tabular}{|c|c|c|}
\hline
\textbf{División} & \textbf{Cociente} &
\textbf{Residuo} \\
\hline
13 / 2 & 6 & 1 \\
6 / 2 & 3 & 0 \\
3 / 2 & 1 & 1 \\
1 / 2 & 0 & 1 \\
\hline
\end{tabular}
\end{center}

\begin{verbatim}
Entrada: 13
Residuos apilados: 1, 0, 1, 1
Salida: 1101
\end{verbatim}

\textbf{Explicación:} Los residuos de las divisiones
se almacenan en la pila. Al retirarlos mediante
\texttt{pop()}, se obtiene el orden correcto del
número binario.
\href{\href{}{}}{}
\section*{Pregunta de cierre}

\textbf{¿Qué característica común tienen los cinco
problemas que hace apropiado utilizar una pila?}

Los cinco problemas utilizan el principio LIFO
(Last In, First Out), que establece que el último
elemento en ingresar es el primero en salir.

Esta característica permite invertir palabras,
recuperar la última acción realizada, verificar
paréntesis balanceados, retirar números en orden
inverso y convertir números decimales a binarios
mediante el almacenamiento de residuos.

\section*{Conclusión}

La realización de estos ejercicios permite comprender
el funcionamiento de las pilas y aplicar sus operaciones
básicas en diferentes problemas de programación.
Se demuestra que el principio LIFO resulta útil cuando
se necesita procesar elementos en orden inverso
al que fueron almacenados.

\end{document}